\section{3D NAND：SLC 电流求和、真实页映射与更新服务}

本文选择二阈值 SLC 的 BL 输入、SL 电流求和路径。在共同 $128\times128$ INT8 任务上，采用文献已经给出尺寸的 NAND 子阵列，把每页的富余位置用于同一逻辑 bit 的等权复制，以接近文献的电流和负载条件。它仍按 R0 逐输入 bit 求值，不采用多位输入加速。由此可得到条件性的读侧范围和完整更新操作计数；现有资料不能给出该 SLC-CIM 状态的完整 program/erase 绝对时间，因此写吞吐和 ridge 保留为可复算函数。该限制不会用另一种状态模式的最快数字补齐。

\subsection{状态、局部实现和证据层级}

\sid{NAND-05} 的实际器件是 16-layer、64 Gb SLC SGVC 芯片。PDF 第1--2页和第3页 Fig.6 给出：在 $V_g=1$ V、$V_{BL}=0.2$ V、$V_{pass}=4.5$ V 条件下，导通电流均值约2 nA、标准差约0.3 nA，关断电流低于0.5 pA。芯片上的器件电流测量支持二阈值操作；论文中约1 $\mu$s 的多位 SA 输出和系统性能属于设计预估。\sid{NAND-04} 使用相关器件数据建立32 WL、64-block 子阵列的 RC/HSPICE 模型，其普通逻辑外围为32 nm，高压驱动另行保留。这里明确区分16层实测器件与32 WL参考模型，不称两者为同一芯片的完整两路实测，也不把相关来源重复计为独立验证。\cite{nand05,nand04}

在所选路径中，data 1 对应低阈值、data 0 对应高阈值；两种状态都需 write-verify 细调（\sid{NAND-04} PDF p.2 II-A）。每次选中一个 WL，将其余 WL 置于 pass 电压，由 BL 施加输入二电平。只有输入与权重同时为1的 NAND 串贡献电流；同一块内的电流在 SL 汇合。原生串、选中/未选中 WL 的偏置、页缓冲编程及高压器件仍需要保留。普通存储的逐 bit 页缓冲判决、商品 $t_R$、多阈值搜索或 read-retry，不是此处 SL 求和的服务周期。CIM 必须另组织 SL 前端、ADC、位权重构和校准。因此既不直接使用普通读时间，也不从普通读时间中凭空拆出一个 cell 延时。

\subsection{固定逻辑状态如何落在原生块与页上}

保留 \sid{NAND-04} PDF pp.2--3 的原生结构：每块 $13824\,BL\times32\,WL\times3\,SSL$，每 subarray 64块、BL长144 $\mu$m。一个选定 WL、SSL 对应的页沿全部13824条 BL 展开，为13824 bit，即1728 Byte。这个页大小由作者所设计的 CIM 结构推导，并非 \sid{NAND-01} 的商品16 KiB页；本文没有把商品 NAND 任意裁成128行的小阵列。独立 SSL 可承载不同权重 bit 的结构见 \sid{NAND-04} Fig.1，写流程仍须满足该页选择粒度。

共同任务有八个二进制权重平面，每轮要求16个输出。因此两个原生子阵列共128块，分别提供 $8\times16$ 个 SL 标量结果；二者共同实现一份状态、共同接收一份输入，由一个更新控制域顺序编程。它们不是两份独立输入服务单元。每块选用一个 SSL；另外两个 SSL 不参与求和。八个输出 WL 组依次服务 $8\times16=128$ 个输出，所以一块实际放置八个有用数据页。输出 $j$、权重位 $k$ 的映射可写为
\[
\text{block}=(k,\,j\bmod16),\qquad
\text{data-WL}=\lfloor j/16\rfloor,\qquad k=0,\ldots,7.
\]

若只驱动128条 BL，最大电流仅256 nA，不能直接引用在大量求和电流下得到的 SL 建立时间。本文因而取 $c=13824/128=108$：将每个逻辑输入及对应权重 bit 等权复制108次，恰好占满物理页。每次仍只处理一个输入 bit，位展开、ADC批数和逻辑字节不变。最大电流由同一2 nA锚点推导为27.648 $\mu$A，与 \sid{NAND-05} 约10 $\mu$A量级、最高约25 $\mu$A的求和场景接近；ADC的参考/增益按108倍电流标定，使码值仍表示 $0\ldots128$ 的近似逻辑部分和。数字阶段不再额外做一次除108运算。

这是使已报告 SL 负载/电流条件较接近本映射的\emph{条件选择}，并非已经验证的最佳编码。在 cell 误差完全独立的理想条件下，满量程电流标准差约为 $0.3\sqrt{13824}$ nA，相对于一个逻辑计数步长216 nA约为0.163 LSB。相关漂移、背面图案依赖和参考误差不会由复制自动消除；这个估计不能证明 INT8 精确计算或最终算法精度。共同合同仍为经校准的 ACIM 近似部分和与24-bit结果容器。源论文7-bit ADC的算法结果不替代本参考名义10-bit、有效约8-bit的外围条件。

数据页共有 $128\times8=1024$ 页，每页只承载128个逻辑权重 bit，即16 Byte的\emph{位片段}，但物理装入为1728 Byte。八个权重平面的对应页合起来才完成128个 INT8 权重。主 resident 事务取完整矩阵，$B_R=16384$ Byte，对应1024个数据页、共1769472 Byte编码数据。物理阵列总规模保持169869312 cell；未用页与物理复制不能增加逻辑 payload。

每块虽有 $32\times3=96$ 个 WL/SSL页位置，固定映射只把其中八个页用于当前逻辑数据，另保留一个或两个 WL 的已知码作校准。其他 WL/SSL 不自动视为可消费的 append 代际。若将未来版本迁往未用 WL，需要另定义地址切换、旧/新版本共存、校准和可持续擦除调度；本分析不借此把擦除分摊到32或96个有用页。

\subsection{Streaming：一次输入的全部求值与输出}

本地宽向量接口接收 $B_S=128$ Byte；输出端交付128个24-bit结果。每个输出 WL 组先完成 WL 建立，再保持该 WL，连续做八个输入 bit 的求值。每轮八个权重平面并行、每平面16个 ADC，128个标量码进入16条重构通道，两拍完成位权合并与输入位累加。补码符号位按负权重构；内部复制、广播、八次输入展开均不重复计算输入 payload。

通过共享 \texttt{acim\_service} 得到
\[
N_E=N_C=N_{rec}=64,\qquad N_{dig}=128,
\qquad N_{scalar}=8192.
\]
八个输出 WL 组需要八次 WL 建立；并非每个 bit 都重新建立 WL，也不能把 WL 开销无限摊薄到外部 workload。\sid{NAND-04} PDF p.3 Table 1给出303 ns WL建立、12 ns BL建立（50\%稀疏条件）以及530--750 ns SL建立；最后一项为大规模输入映射下的 RC/HSPICE估计，SL负载约16 pF。取530、640、750 ns分别配共同短、参考、长外围；640 ns只是区间中点。公共输入时隙包含低压捕获/公共复位，BL与SL项包含本映射专用建立；128个 ADC及数字通道保持 R0 数量。

令 $t_x\ge0$ 为每次求值仍需额外计入的专用复位、跨阻建立或恢复时间。完整步骤表达为
\begin{equation}
\begin{split}
\Delta_S={}&8t_{WL}+64(T_I+t_{BL}+t_{SL}+t_x)\\
 &+64T_A+128T_D+2T_D.
\end{split}
\label{eq:nand-stream}
\end{equation}
最后两拍是输入捕获与结果提交。所有阶段顺序执行，无跨向量流水或异步 plane 并行，也不使用 SSD或die接口带宽。若专用前端/复位能在所列预算内完成，取 $t_x=0$ 得表~\ref{tab:read}；额外阶段有正时间时由式~\ref{eq:nand-stream}直接补入，不能漏计。

% Generated by scripts/check_nand.py --emit.
\begin{table}[htbp]\centering\small
\begin{tabular}{@{}lrrrrr@{}}\toprule
情景 & $t_{SL}$ (ns) & $\Delta_S$ ($\mu$s) & $\rho$ (MB/s) & $F_{app}$ ($\mu$s) & $F_{rw,1}$ ($\mu$s)\\\midrule
短时隙 & 530 & 38.140 & 3.356 & 225.28 & 253.44\\
参考 & 640 & 46.402 & 2.759 & 563.2 & 633.6\\
长时隙 & 750 & 56.332 & 2.272 & 1126.4 & 1267.2\\
\bottomrule\end{tabular}
\caption{条件读侧结果与整矩阵写数据装入/控制开销。$F_{app}$ 为1024数据页；$F_{rw,1}$ 为1024数据页加128校准页。MB 为$10^6$ Byte。SL中点640 ns为插值；不是新增实测点。}
\label{tab:read}\end{table}


该条件范围为 $\rho\approx2.27$--$3.36$ MB/s（约 $0.00227$--$0.00336$ GB/s）。它要求原生几何、复制后电流、约16 pF SL负载与所采用530--750 ns预算相容，并要求新的28 nm前端达到规定误差；12 ns BL值也保留其原负载条件。复制只能改善与来源的接近程度，不能证明这些条件已被硅验证。若条件不成立，读侧只能保留式~\ref{eq:nand-stream}而非把此数值称为保证范围。\sid{NAND-05} 的约1 $\mu$s多位输出仅作量级交叉检查，不与上述时间重复相加。

\subsection{Resident：完整编程、擦除和校准维护}

从本地更新命令开始，到整矩阵全部 INT8 权重可用于上述求值结束，计 $B_R=16384$ Byte。128-bit接口是编码数据装入宽度，不是可同时编程128 cell的断言。每物理页需 $13824/128=108$ 个数据拍；页命令/数据分开，准备和完成共两拍。利用共享 \texttt{front\_ns} 得
\[
F_{page}=110T_D.
\]
本参考提供至少一页1728 Byte的页缓冲；同一更新域逐页装入、执行完整program并确认完成。阵列内部允许一页目标 cell并行处理，但1024页或128块不能未经证据同时执行。

定义 $P$ 为\emph{一个完整 SLC-CIM 页program}的时间，含高压泵/驱动建立、全部program脉冲、两阈值目标所需verify及重试、恢复与到可计算状态的终点；$E$ 为一个完整block erase，含erase-verify和恢复。源文约20 V编程驱动不能用0.9 V低压写拍替代。若将来获得完整周期，内部verify已包含则只计一次；不能再逐轮叠加，亦不能用一个脉冲代替完整 $P$。

\textbf{预擦除append窗口。} 八个数据WL已有预擦除空间，校准WL已有已知码。按固定映射首次写入本矩阵后，该窗口结束；不默认后台有无限擦除空间。额外校准时间记为整矩阵事务的 $C_{app}$。只有新权重布局及背面图案仍满足已有校准时才能取零，否则必须计入：
\[
\Delta_{R,app}=1024(F_{page}+P)+C_{app}.
\]

\textbf{同地址持续整矩阵重写。} 任意新权重不能保证只向一个阈值方向变化，因此先顺序擦除128块，再重建数据与校准页。\sid{NAND-05} PDF p.2 II-(e)、p.4 Fig.12建议每块用一或两个已知码WL作SA校准。擦除会同时销毁这些页，它们不是免费保存的参考状态。主情景每块保留一个校准页，重写需1024个数据页加128个校准页program；若需要两个校准页则加256页。令 $q=1$ 或2，$C_{rw}$ 为\emph{整个矩阵事务}在这些页编程完成后的额外感测、增益/参考配置、必要重校准总时间，非每块时间，则
\begin{equation}
\Delta_{R,rw}=(1024+128q)(F_{page}+P)+128E+C_{rw}.
\label{eq:nand-write}
\end{equation}
校准页采用同一二阈值program引擎和相同完整 $P$；若其终点不同，应按实际不同周期替换。这些内部写入和校准都不增加 $B_R$。维护调度不得偷偷与求值重叠。冷启动也需要建立校准页和有效校准，不能直接使用预校准append初始条件。

% Generated by scripts/check_nand.py --emit.
\begin{table}[htbp]\centering\small
\begin{tabular}{@{}lrrrl@{}}\toprule
事务 & program页数 & erase块数 & $F$ ($\mu$s) & $\Delta_R$ ($\mu$s)\\\midrule
预擦除append & 1024 & 0 & 563.2 & $563.2+1024P+C$\\
重写，1校准WL & 1152 & 128 & 633.6 & $633.6+1152P+128E+C$\\
重写，2校准WL & 1280 & 128 & 704 & $704+1280P+128E+C$\\
\bottomrule\end{tabular}
\caption{参考外围下的完整事务系数。每行$B_R=16384$ Byte；$P,E,C$均以$\mu$s计，分别为完整SLC-CIM页program、块erase、整事务额外校准时间。三者未被本地证据定值。}
\label{tab:coeff}\end{table}


擦除摊销对应每块八个有用数据页，故把持续整矩阵更新换算为每数据页的等效占用时，擦除项是 $E/8$，主情景program与装入项是 $\tfrac98(P+F_{page})$。每等效页的16 Byte只是核算单位，不能说一个位平面页就完成了16个 INT8 权重。完整矩阵事务和聚合后系数一致。任意局部一个输出的原位更新若破坏同块其他七个输出，必须另计保留状态的读出、重编程及可计算性；本结果仅覆盖明确的全矩阵替换。

\subsection{证据可以给出的范围、缺口和主导因素}

本地互补资料已查完。\sid{NAND-01} 给的是MLC/TLC两列商品 $t_{PROG}$、$t_{BERS}$ 与读时间；\sid{NAND-03} 是QLC的1.3 ms页program和6-plane、75 MB/s die吞吐；\sid{NAND-02} 给QLC芯片SLC-cache burst约11\%的相对改善，未给其SLC绝对program或erase。burst图还显示首/末页必须支付泵建立/最终放电。这些证据能解释操作结构，却不能构成所选二阈值电流细调状态的绝对上界或下界。故不把任何MLC/TLC/QLC时序代入 $P,E$，不采用相对burst收益，也不填造一个SLC典型值。\cite{nand01,nand03,nand02}

三个公共外围情景给出明确的装入/控制开销，未知 $P,E,C$ 的系数均已计算。以参考外围、一个校准WL为例，用 $\mu$s统一时间单位，记 $C=C_{rw}$：
\begin{align}
\tau_{app}&=\frac{16384}{563.2+1024P+C_{app}}\ \mathrm{MB/s},\\
\tau_{rw}&=\frac{16384}{633.6+1152P+128E+C}\ \mathrm{MB/s},\\
\RI^*_{app}&=\frac{4.4+8P+C_{app}/128}{46.402},\\
\RI^*_{rw}&=\frac{4.95+9P+E+C/128}{46.402}.
\label{eq:nand-ridge}
\end{align}
这里采用共享关系
\[
\rho\approx B_S/\Delta_S,\quad \tau\approx B_R/\Delta_R,\quad
\RI^*\approx (B_S/B_R)(\Delta_R/\Delta_S),\qquad B_S/B_R=1/128.
\]
$P,E,C$未定，因而不能报告有文献支持的有限数值 $\tau$或 $\RI^*$范围；JSON相应字段为 \texttt{null}，并保留完整线性系数。两校准WL情景将program系数从1152改为1280，参考装入/控制项由633.6改为704 $\mu$s；逻辑容量和求值吞吐不变。

% Generated by scripts/check_nand.py --emit.
\begin{table}[htbp]\centering\small
\begin{tabular}{@{}lrrr@{}}\toprule
外围情景 & append的$P$阈值 & 重写1校准WL的$P$阈值 & 重写1校准WL的$E$截距\\\midrule
短时隙 & 4.548 & 4.018 & 36.160\\
参考 & 5.250 & 4.606 & 41.452\\
长时隙 & 5.941 & 5.159 & 46.432\\
\bottomrule\end{tabular}
\caption{由$\RI^*=1$反推的时间预算（$\mu$s）。$P$阈值设$E=C=0$；$E$截距设$P=C=0$。零值只用于求直线截距，不是物理情景或器件时间。}
\label{tab:threshold}\end{table}


表~\ref{tab:threshold}给出另一种有数值的结论：它把 $\RI^*=1$ 对应的可容许program/erase预算反算出来，而不是假定器件具备这些时间。参考主情景满足平衡线
\[
9P+E+C/128=41.452\ \mu\mathrm{s}.
\]
任何实测 $P,E,C$代入后即可判定ridge位于该阈值哪侧。这一预算对写验终点和物理页利用率极为敏感；它不能从一次脉冲或die吞吐直接判断。

读侧主项是64次SL建立，三个情景分别贡献33920、40960、48000 ns，约占总服务时间89\%、88\%、85\%；八次WL建立固定2424 ns。共同低压外围改变会影响结果，但尚未闭合的SLC完整写验、擦除及校准更根本地限制写侧结论。每页只有16 Byte逻辑位片段而装入1728 Byte，使本固定小逻辑状态难以直接享有大容量NAND的全部页吞吐。此处不是NAND材料的固有吞吐排名。

\subsection{对共享方法的反馈与适用边界}

该案例遵循R1的共同逻辑任务、近似精度合同和28 nm外围；R2记录原生block/SL输出及WL分组；R3记录108次等权复制的资源与适配条件；R4将完整周期与组件时序分开；R5保持真实页、有效payload、单更新域和校准维护。机器输入中逐项保存R0值、变更、证据和两路影响。

本案例有三项接入注意事项，均可由现有顺序组合与例外规则处理，无需改变共享理论。第一，WL建立的调用次数与BL/SL求值次数不同，单一每轮 $t_m$不足以表达；本脚本导入共享接口后，以同一顺序操作模板追加八次WL块，未复制另一套吞吐公式。第二，跨位平面的页并不是完整INT8更新事务，必须先合并全部必要页，再计算有用字节。第三，erase会破坏校准参考，而未用页与校准页不等于有用更新页，不能把原生page/block数直接放进摊销分母。

所有结果均为单服务域、顺序执行的参考估算。有效保持、读扰与温度漂移若要求周期校准或重编程，内部维护占用还须加入对应服务窗口，不能记成新的workload更新payload；这里只在新事务中显式保留了必要校准时间变量。耐久度另属使用可行性，现有SLC-CIM资料不足以给出统一寿命上限，也不将商品MLC/TLC寿命除成年写入率混入瞬时 $\tau$。本节不生成最终Table I、不分析workload，也不声称两路独立能力可以同时达到。
